\documentclass[aps,prb,reprint,amsmath,amssymb,superscriptaddress,longbibliography,nofootinbib,floatfix]{revtex4-2}
\usepackage[T1]{fontenc}
\usepackage{lmodern}
\usepackage{graphicx}
\usepackage{bm}
\usepackage{microtype}
\usepackage[hidelinks]{hyperref}
\graphicspath{{../../results/novelty_20261008/four_cell_window_audit/}{../../results/novelty_20261008/figures/}{../../results/novelty_20261004/rank_figures/}{../../results/novelty_20261004/rank_controls_figure/}{../../results/novelty_20261004/evidence_figures/}{../figures/}}
\newcommand{\Dloc}{\Delta_{\mathrm{loc}}}
\newcommand{\Wlow}{W_{\mathrm{low}}}
\newcommand{\Cmin}{C_{\mathrm{min}}}
\begin{document}
\raggedbottom
\title{Thermal stability of retuned interference zeros in a Heisenberg pentamer chain}
\author{Author Name}
\email{corresponding.author@example.edu}
\affiliation{Department and Institution, City, Postal Code, Country}
\date{\today}
\begin{abstract}
Can exchange tuning make a low-energy spin spectrum dark for every
thermally occupied initial state? We address this question in a weakly
coupled isotropic Heisenberg pentamer chain. For an isolated local
doublet, the entire first-order probe can be canceled on a uniform
reflection-symmetric periodic chain, leaving weight of order $\beta^4$
or smaller at fixed $T/J_{\rm eff}$ as the intercell exchange $\beta$
tends to zero. A local quartet instead produces dipole and composite
vectors with the same symmetry labels but different dynamical
matrix-element ratios. At the specified interference root, we prove
that their leading positive-frequency thermal response cannot be
canceled for any finite even periodic chain with at least four cells.
The obstruction holds at every positive scaled temperature, even with
independent first-order cell-local SU(2)-invariant exchange controls.
An exact spin-dependent second moment quantifies it, and locality gives
a size-independent finite-frequency weight bound in the effective
uniform-trace limit. The leading probe nevertheless annihilates a
four-cell ground state; its first thermal response has a calculable activation coefficient
and line-weight ratio $9:8$. Complete low-band spectra of twenty
microscopic spins verify these lines and expose a fourth-order
background at low temperature. Finite momentum acceptance supplies a
separate background. These results identify a limit on local exchange
control of thermally activated scattering, rather than a new symmetry
selection rule or a resolution-independent experimental diagnostic.
\end{abstract}
\maketitle
\section{Introduction}

Weak intensity in a spin spectrum need not mean an absence of excitations.
In a system with several magnetic sites per unit cell, destructive interference
can suppress the matrix element of the physical probe while the collective
mode remains present. Cluster interference and symmetry-constrained
momentum dependences are established tools in magnetic
scattering~\cite{waldmann2003interference,haraldsen2011pentamer,tabrizi2021selection,tabrizi2021anisotropic}.
Effective chains of cluster spins likewise provide a familiar route to
low-energy modes and folded spectra, including trimer and tetramer
chains~\cite{cheng2022trimer,li2026tetramer}.

The projected form factor is only the leading approximation to a scattering
amplitude. Perturbative corrections to dimer structure factors already show
that interactions can modify this momentum dependence~\cite{mueller2003structure}.
Waldmann and G\"udel constructed effective scattering operators beyond simple
spin projection and analyzed tensor corrections caused by anisotropy-induced
mixing of cluster multiplets~\cite{waldmann2005manyspin}.
More recently, systematic extinctions have been distinguished from zeros that
depend on Hamiltonian parameters~\cite{tabrizi2026extinctions}.
Neither a form-factor zero nor the use of a corrected effective observable is
therefore the question addressed here.

The failure of symmetry to fix an intensity is also established.
For isotropic clusters, Ghassemi Tabrizi derives a universal momentum
dependence when a unique transition-operator irrep mediates a transition
and occurs once in the site-spin operator representation.
When the irrep occurs repeatedly, Eq.~(16) of
Ref.~\onlinecite{tabrizi2021selection} instead retains independent reduced
matrix elements and their interference. Symmetry therefore need not fix
their ratio. This does not decide whether the ratios for all thermally
accessible transitions can be adjusted together by physical exchange
controls. That simultaneous cancellation is the issue addressed below.

Temperature makes the interpretation of a dark point more demanding.
Effective-observable vertex corrections already describe conditional
excitations and thermal changes of spectral weight in spin
ladders~\cite{exius2010vertex}. Thermally activated multiparticle
transitions and spectral-weight transfer have also been analyzed in
frustrated ladders~\cite{honecker2016thermal}. Thermal lifting of an
accidental intensity cancellation has also been observed for interfering
nuclear and magnetovibrational phonon amplitudes~\cite{raymond2018interference}.
The appearance of additional lines on heating is therefore not itself a
new mechanism. The question
here is whether an exchange parameter can be retuned to cancel all those
transitions together, rather than only the ground-state response.

The Heisenberg pentamer chain provides doublet and quartet local ground
multiplets without changing the microscopic spin length or bond topology.
We show that their retuned zeros have different thermal stability. A
doublet permits cancellation of the entire first-order probe, whereas a
quartet produces a composite vector whose bond-spin channels cannot all be
canceled by the same real retuning. A local commutator argument establishes
this obstruction at nonzero scaled temperature for every finite even
periodic chain in the specified effective model. An exact second-moment
sum rule extends the algebraic obstruction to other local spins,
conditional on the same dipole--composite probe structure. Combined with
locality, it gives a size-independent lower bound on finite-frequency
weight in the uniform-trace limit. Allowing independent first-order
SU(2)-invariant retunings within each cell does not lower the optimal
response below uniform retuning: additional operator momenta contribute
nonnegative spectral weight. A four-cell ring is dark
at zero temperature but acquires an exactly calculable activated response.
Microscopic spectra test the weak-coupling powers and the four-cell activated
line groups. Finite-coupling backgrounds and finite momentum
acceptance limit the use of this distinction
as an experimental diagnostic.

\section{Model and low-energy observable}

Each cell contains five spin-$1/2$ sites at unit spacing. With $J_1=1$,
$\alpha=J_2/J_1$, and $\beta=J_3/J_1>0$, the Hamiltonian is
\begin{align}
 H_0&=\sum_{n=0}^{N-1}\big[
 \alpha(\bm S_{n,1}\!\cdot\!\bm S_{n,2}
       +\bm S_{n,4}\!\cdot\!\bm S_{n,5})\notag\\
 &\hspace{29mm}+\bm S_{n,2}\!\cdot\!\bm S_{n,3}
       +\bm S_{n,3}\!\cdot\!\bm S_{n,4}\big],\\
 H&=H_0+\beta V,\qquad
 V=\sum_{\langle n,n+1\rangle}\bm S_{n,5}\!\cdot\!\bm S_{n+1,1}.
 \label{eq:rank-ham}
\end{align}
The bond sum is specified as open or periodic in each calculation, and the
microscopic length is $L=5N$. We assume that the isolated cell has a single
ground multiplet of spin $s$, separated from all other multiplets by
$\Dloc>0$. The numerical examples below have $s=1/2$ or $3/2$.

Rotational symmetry within this irrep fixes the local projection,
\begin{equation}
 P_n\bm S_{n,a}P_n=c_a\bm T_n,\qquad
 c_a=c_{6-a},\qquad \sum_a c_a=1.
 \label{eq:rank-projection}
\end{equation}
Here $\bm T_n$ is a spin-$s$ operator. The leading Hamiltonian is an
antiferromagnetic effective chain with exchange $\beta c_1^2$.
This standard reduction does not by itself determine a suppressed intensity.

It is convenient to center the intra-cell Fourier phase:
\begin{align}
 B_n^z(q)&=\sum_{a=1}^5 e^{iq(a-3)}S_{n,a}^z,\\
 O_q&=\frac{1}{\sqrt{5N}}\sum_n e^{iQn}B_n^z(q),
 \quad Q=5q\pmod{2\pi}.
\end{align}
Changing the coordinate origin adds only a global phase to $O_q$ and does not
change a spectral weight. Equation~\eqref{eq:rank-projection} gives
\begin{align}
 P O_qP&=A(q,\alpha)D_Q,\notag\\
 A(q,\alpha)&=\sum_a c_a\cos[q(a-3)],\\
 D_Q&=\frac{1}{\sqrt{5N}}\sum_n e^{iQn}T_n^z.
 \label{eq:rank-amplitude}
\end{align}
The normalization of $D_Q$ is inherited from the microscopic probe, not
chosen independently for each curve.

Our observable is the integrated inelastic low-band weight
\begin{equation}
 \Wlow=\frac{1}{n_G}\sum_{r=1}^{n_G}
 \sum_{0<E_m-E_G<\Dloc/2}|\langle m|O_q|G_r\rangle|^2.
 \label{eq:rank-weight}
\end{equation}
All transitions within the ground eigenspace are excluded.
The gap $\Dloc$ in this cutoff and in $g=\beta/\Dloc$ is evaluated at the
unshifted local root. For each tested finite system we check that the cutoff
contains the complete retained low band. This is a fixed-size,
weak-coupling construction, not a claim of a uniformly isolated many-body
band as $N$ tends to infinity.

\section{What survives a retuned zero?}

\subsection{Scalar and quadrupolar local factors}

Let $P=\prod_n P_n$ and $R=(1-P)/(H_0-NE_5)$, where $E_5$ is the isolated
ground-multiplet energy. At a root $A(q,\alpha_*)=0$, the leading effective
probe is
\begin{align}
 (O_q)_{\mathrm{eff}}&=\beta K_q+O(\beta^2),\notag\\
 K_q&=-P O_q R V P-P V R O_q P.
 \label{eq:rank-probe}
\end{align}
The two terms must be evaluated separately when $O_q$ is non-Hermitian.
They include virtual local excitations absent from the first-order
Hamiltonian.

A link can contribute only through an intermediate excitation on the cell
acted on by the probe: the local probe cannot de-excite both cells of a link.
For a Hermitian component of $B^z$, the relevant local tensor is
\begin{equation}
 M^{z\mu}=P_n(B^z R_n S_{\mathrm{end}}^\mu
                   +S_{\mathrm{end}}^\mu R_n B^z)P_n.
 \label{eq:rank-local-tensor}
\end{equation}
It is Hermitian and even under time reversal. Coupling two vectors permits
ranks zero, one, and two, but a Hermitian rank-one tensor within a single
spin irrep is odd under time reversal. Only the scalar and quadrupolar parts
remain. The real and imaginary Hermitian components of $B(q)$ can be treated
separately and recombined.

Define
\begin{equation}
 Q_n^{z\mu}=\tfrac12\{T_n^z,T_n^\mu\}
        -\tfrac13 s(s+1)\delta_{z\mu}I.
\end{equation}
On one orientation of a link the correction consequently has the form
\begin{equation}
 \delta B_n^z=a\,T_{n+1}^z
       +b\sum_\mu Q_n^{z\mu}T_{n+1}^\mu,
 \label{eq:rank-link}
\end{equation}
with a corresponding term for the opposite endpoint.
The second term is a vector under simultaneous spin rotations. It does not
represent a direct rank-two transition or a new particle.

\subsection{Cancellation for a local doublet}

In a single spin-$1/2$ irrep, $Q_n^{z\mu}=0$ identically. Equivalently every
Hermitian local matrix is a linear combination of $I$ and $\bm T$; time
reversal leaves only $I$ in Eq.~\eqref{eq:rank-local-tensor}.
The entire first-order correction is thus a sum of neighboring one-spin
operators. On a uniform periodic chain at a commensurate momentum, translation
reduces this sum to $K_q=\gamma(q)D_Q$. Cell-centered reflection and real
exchange give $a_R=a_L^*$ for the two link orientations. Consequently
$\gamma=a_L e^{-iQ}+a_R e^{iQ}=2\operatorname{Re}(a_L e^{-iQ})$
is real in the same phase convention as $A$.

Suppose the local zero is simple, with
$A'=\partial_\alpha A(q,\alpha_*)\ne0$. A physical exchange retuning
$\alpha=\alpha_*+\kappa\beta$ gives
\begin{equation}
 (O_q)_{\mathrm{eff}}
   =\beta[\gamma+\kappa A']D_Q+O(\beta^2).
\end{equation}
Choosing $\kappa=-\gamma/A'$ cancels this operator throughout the retained
Hilbert space, not only on a selected transition. Hence
\begin{equation}
 \Wlow=O(\beta^4)\quad\text{or smaller}.
 \label{eq:rank-doublet}
\end{equation}
The conclusion requires one isolated irrep without additional local
multiplicity, SU(2) and time-reversal symmetry, bilinear links, and the
periodic/reflection conditions just stated. It does not apply unchanged to
open endpoints. Nor does it establish a nonzero fourth-order coefficient
after arbitrary higher-order retuning.

\subsection{Residual for a quartet}

For $s=3/2$ the quadrupole in Eq.~\eqref{eq:rank-link} is nonzero, and a
single retuning need not remove its source. Define $d=D_Q|G\rangle$ and
$k=K_q|G\rangle$, with the entire ground subspace projected out.
For degenerate ground states these vectors denote a direct sum, with norm
averaged over the ensemble used in Eq.~\eqref{eq:rank-weight}.
The leading weight after a real retuning is
\begin{equation}
 \Wlow=\beta^2\|k+\kappa A'd\|^2+O(\beta^3).
\end{equation}
Elementary least squares gives
\begin{align}
 \kappa_{\mathrm{opt}}
   &=-\frac{\operatorname{Re}\langle d,k\rangle}{A'\|d\|^2},\\
 \Cmin&=\|k\|^2-
       \frac{[\operatorname{Re}\langle d,k\rangle]^2}{\|d\|^2}.
 \label{eq:rank-minimum}
\end{align}
The overlap is real to numerical precision in the examples below.
A nonzero $\Cmin$ measures the part that cannot be removed by the leading
one-parameter retuning. The algebraic identity is standard; the physical
question is whether the corrected source in this model is noncollinear with
$d$. An allowed tensor alone does not settle that question.

At $q=\pi$ the quartet branch provides a simple analytic example:
\begin{align}
 \alpha_*&=-\frac{2+\sqrt7}{6},\notag\\
 c_1=c_5&=\frac{25+4\sqrt7}{108},\quad
 c_2=c_4=\frac14,\notag\\
 c_3&=\frac{1-2\sqrt7}{27}.
 \label{eq:rank-exact-root}
\end{align}
These coefficients obey $c_1-c_2+c_3-c_4+c_5=0$.
Here $\Dloc=0.16622155$, $c_1^2=0.10855198$, and
$A'=0.10026088$. Exact local resolvent evaluation gives
\begin{align}
 a&=\frac{4500505-2595746\sqrt7}{19525536},\notag\\
 b&=\frac{1082642\sqrt7-2261005}{4881384}>0.
 \label{eq:exact-local-coefficients}
\end{align}
Thus $a=-0.12123577$ and $b=0.12361176$ in
Eq.~\eqref{eq:rank-link}: the required composite coefficient is exactly
nonzero, rather than a residual inferred from a numerical tensor fit.
The construction is described in Appendix~\ref{app:exact-local-check}.

\subsection{A bond-spin interpretation}

The composite vector also has a direct interpretation in terms of neighboring
total-spin channels. For two spin-$s$ sites, put
$h_{ij}=\bm T_i\cdot\bm T_j$, $D_{ij}=T_i^z-T_j^z$, and
$X_{ij}=\sum_\mu(Q_i^{z\mu}T_j^\mu-T_i^\mu Q_j^{z\mu})$.
Expanding the quadrupoles gives the exact identity
\begin{equation}
 X_{ij}=\tfrac12\{h_{ij},D_{ij}\}+\tfrac13s(s+1)D_{ij}.
 \label{eq:bond-identity}
\end{equation}
Within bond total spin $J$, $h_{ij}$ has eigenvalue
$[J(J+1)-2s(s+1)]/2$. The exchange-odd vector $D_{ij}$ connects only
adjacent bond-spin sectors. If $\Pi_J$ denotes their projector, then
\begin{equation}
 \Pi_J X_{ij}\Pi_{J-1}=
 \left[\frac{J^2}{2}-\frac{2s(s+1)}{3}\right]
 \Pi_J D_{ij}\Pi_{J-1}.
 \label{eq:bond-channel}
\end{equation}
For a doublet there is only $J=1\leftrightarrow0$, and this coefficient
vanishes. For a quartet, the $1\leftrightarrow0$, $2\leftrightarrow1$,
and $3\leftrightarrow2$ coefficients are $-2,-1/2,2$.
The link probe $-aD_{ij}+bX_{ij}$ therefore renormalizes these channels
differently, whereas exchange retuning adds the same multiple of $D_{ij}$.
One adjustment cannot set all three link-operator coefficients to zero.
Whether those independent channels leave weight on a particular chain ground
state still requires Eq.~\eqref{eq:rank-minimum}: interference between links
and restricted ground-state matrix elements can remove the residual.

The two types of zero are not confined to the single comparison in
Table~\ref{tab:rank-roots}. Local diagonalization found simple roots at each
of 16 sampled values of $Q/\pi$ from $0.55$ to $1$, using
$q=(Q+4\pi)/5$ for the quartet and $q=(Q+2\pi)/5$ for the doublet
[Fig.~\ref{fig:rank-controls}(a)]. All sampled local gaps were positive,
with minimum $0.04263$. The quartet coefficient $b$ was nonzero throughout
this sample, whereas the doublet quadrupole vanished identically.
This local calculation identifies a family of candidate interference settings;
it is neither a phase diagram nor evidence that the many-body residual is
nonzero at every momentum in the interval.

\section{Thermal stability of a retuned zero}
\label{sec:thermal}

\subsection{An operator zero and a ground-state zero}

A vanishing ground-state source does not establish that a probe vanishes as
an operator. Temperature distinguishes these possibilities by populating
additional initial states. We keep $\tau=T/J_{\rm eff}$ fixed as
$\beta\to0$, where $J_{\rm eff}=\beta c_1(\alpha_*)^2$ and $k_B=1$.
Thus the thermal ensemble remains inside the local ground multiplets.
Put $h=\sum_i\bm T_i\cdot\bm T_{i+1}$ and let $\epsilon_i$ be its
dimensionless eigenvalues. For the Hermitian leading $Q=\pi$ probe,
write $D=D_\pi$, $K=K_q$, and define
\begin{align}
 R_\kappa&=K+\kappa A'D,\notag\\
 C_\kappa(\tau)&=\sum_{\epsilon_f>\epsilon_i}
 p_i\left|\langle f|R_\kappa|i\rangle\right|^2,
 \qquad p_i=\frac{e^{-\epsilon_i/\tau}}{Z}.
 \label{eq:thermal-weight}
\end{align}
This coefficient describes the regular positive-frequency response on the
$J_{\rm eff}$ scale. At fixed finite size, choose $\eta>0$ smaller than
the smallest nonzero level spacing of $h$, and let $W_{+,\eta}$ include
only low-band transitions with $\omega>\eta J_{\rm eff}$. Then
$W_{+,\eta}=\beta^2C_\kappa+O(\beta^3)$. This qualification matters:
higher-order terms can split degenerate levels of $h$ and move weight
from zero frequency to $\omega=O(\beta^2)$. That quasielastic weight is
not automatically included in $C_\kappa$, even though it contributes to
a microscopic integral over every $\omega>0$.
All equal-energy blocks, including off-diagonal elements within degenerate
subspaces, are excluded. A real exchange retuning is optimized separately
at each temperature, so a mere displacement of the zero cannot account for
a nonzero minimum.

At any $\tau>0$ every state of a finite effective chain has positive
Boltzmann weight. Nonnegativity of the summands and Hermiticity imply
\begin{equation}
 C_\kappa(\tau)=0\quad\Longleftrightarrow\quad[h,R_\kappa]=0.
 \label{eq:thermal-criterion}
\end{equation}
The right-hand side allows an operator acting within equal-energy blocks;
it need not be the zero operator. This standard positivity observation is
useful here because it replaces a temperature-by-temperature cancellation
test with a local operator calculation. For a doublet, the stronger identity
$K=\gamma D$ already proved above cancels $R_\kappa$ throughout the
retained space. Consequently $W_+=O(\beta^4)$ or smaller at every fixed
$\tau$ under the same periodic and reflection assumptions.

\subsection{A representation-dependent obstruction}

For the quartet root at $q=\pi$, define the unnormalized alternating spin
$\widetilde D=\sum_i(-1)^iT_i^z$ and
$Y=\sum_i(-1)^iX_{i,i+1}$, using Eq.~\eqref{eq:bond-identity}. Then
\begin{equation}
 R_\kappa=\frac{bY+(-2a+\kappa A')\widetilde D}{\sqrt{5N}}.
 \label{eq:thermal-source}
\end{equation}
Both $\widetilde D$ and $Y$ are spin vectors with momentum $\pi$,
odd under time reversal and even under the site-permutation reflection
$i\mapsto-i$. For $Y$, reflection reverses the bond orientation and
shifts its alternating sign, leaving their product unchanged.
The two contributions therefore share these symmetry labels. The issue is
not whether a transition is symmetry-allowed, but whether one dipole
coefficient can cancel all its dynamical amplitudes across different
initial and final states.
This is consistent with the independent reduced matrix elements in the
multiplicity case of Ref.~\onlinecite{tabrizi2021selection}, Eq.~(16).
Their independence as operator structures alone is not a proof of
noncancellation: the inelastic matrix elements could still be
proportional, with their difference acting only inside equal-energy
blocks. The commutator minimum below excludes precisely this possibility
for the specified effective chain.
Different bond-spin channels obstruct cancellation of its commutator with
$h$. The spectral-moment approach is standard in spin-chain
dynamics~\cite{mueller1982sumrules}; here it gives an exact retuning
minimum for this dipole--composite probe. More generally, take a spin-$s$
chain with the same uniform nearest-neighbor
$h$ and probe $R=(bY+u\widetilde D)/\sqrt{\ell N}$, where
$\ell>0$ fixes the probe normalization ($\ell=5$ for the pentamer).
Let $d=2s+1$ and $c=s(s+1)$. For every finite even $N\ge4$, the
uniformly traced, two-sided second frequency moment is exactly
\begin{align}
 M_2(u)&=\frac{1}{d^N}\operatorname{Tr}[h,R]^\dagger[h,R]\notag\\
 &=M_{2,\min}+\frac{(d^2-1)^2}{18\ell}
       \left[u+\frac{b(d^2-4)}{12}\right]^2,\\
 M_{2,\min}&=\frac{b^2c^2(4c-3)(86c-27)}{810\ell}.
 \label{eq:general-moment}
\end{align}
Appendix~\ref{app:general-moment} derives this identity from local spin
traces and the orthogonality of distinct operator supports. It is not a
fit to a finite set of spins or chain lengths. Independent full-chain
matrices for $s=1/2,1,3/2$ at $N=4,6$ reproduce the identity.

For $s=1/2$ the quadrupole vanishes and the minimum is zero. For
$s\ge1$ and $b\ne0$, it is strictly positive. The conclusion is
conditional on this probe form and on retuning only its dipole
coefficient: a higher local spin alone does not guarantee a nonzero $b$
in an arbitrary microscopic cluster. For the pentamer quartet the exact
minimum is $(197/16)b^2$. Combining this positive minimum with
Eq.~\eqref{eq:thermal-criterion} proves $C_{\min}(\tau)>0$ for every
finite even $N\ge4$ and every $\tau>0$ in the leading quartet model.
A ground-state cancellation therefore cannot be continued into complete
thermal cancellation by the same one-parameter retuning.
The moment optimum $u=-b$ differs from the four-cell ground-state optimum
$u=3b$; minimizing one observable need not minimize another.

This trace includes positive and negative frequencies; its positive-only
uniform-trace value is half as large. It refers to the retained effective
space, not the full microscopic infinite-temperature spectrum. In the
fixed-size weak-coupling limit it describes the ordered regime
$J_{\rm eff}\ll T\ll\Dloc$. At the four-cell zero-temperature retuning,
the prediction is $M_2=(837/16)b^2=0.79932812$. Directly summing the
twenty-spin low-band transitions, with uniform initial weights and
division by $\beta^2J_{\rm eff}^2$, gives $0.79967945$ at
$g=0.000625$, a $0.044\%$ deviation; the deviation increases to
$0.088\%$ and $0.177\%$ at the next two couplings. No frequency window
is used for this check: the frequency-squared factor suppresses the
quasielastic contribution in the leading limit. This does not remove the
separate finite-momentum limitation or by itself yield an integrated-weight
lower bound.

Locality supplies a stronger conclusion specifically in the uniform-trace
limit. Let $C^{\rm tr}_{>\nu}$ denote the positive-frequency weight of
$R$ with dimensionless frequency above $\nu$, using equal initial-state
weights $1/d^N$. Appendix~\ref{app:weight-bound} proves, for the pentamer
quartet and every real $u$,
\begin{equation}
 C^{\rm tr}_{>\sqrt{45}/4}\ge\frac{197}{1036800}b^2>0
 \qquad(N\ge4\text{ even}).
 \label{eq:quartet-window-bound}
\end{equation}
This uses a fourth-moment upper bound as well as
Eq.~\eqref{eq:general-moment}. It excludes a sequence of larger even
rings whose optimized uniform-trace response vanishes or moves entirely
to arbitrarily small frequencies. The coefficient is conservative, not
a predicted intensity. The statement is within the leading effective
model with the uniform-trace limit taken before increasing $N$; it does
not establish a lower bound at arbitrary fixed $T/J_{\rm eff}$ or
uniform validity of the microscopic weak-coupling expansion with size.

\subsection{Independent local controls do not evade the obstruction}

The preceding result might appear to reflect an insufficient number of
control parameters. At the quartet root $q=\pi$, however, any bounded
SU(2)-invariant perturbation of order $\beta$ confined to an individual
cell changes its projected vector probe only by a multiple of $T_i^z$.
The isolated cell contains a single spin irrep; its projected scalar
perturbation shifts all components equally. The leading uniform $h$ and
the coefficient $b$ are therefore unchanged. Even independent local
controls in every cell can only enlarge the leading probe to
\begin{equation}
 R_{\{u_i\}}=\frac{bY+\sum_i(-1)^iu_iT_i^z}{\sqrt{\ell N}},
 \qquad u_i\in\mathbb R.
 \label{eq:local-control-probe}
\end{equation}
This statement concerns infinitesimal local retunings, not arbitrary
changes to the leading intercell exchange or to the local spin multiplet.

Translation symmetry separates the $Q=\pi$ component, whose dipole
coefficient is $\bar u=N^{-1}\sum_i u_i$, from all other operator momenta.
Their thermal spectral weights add without interference. For either the
positive-frequency weight or any fixed positive-frequency window,
\begin{equation}
 F(R_{\{u_i\}})=F(R_{\bar u})+
 F(R_{\{u_i\}}-R_{\bar u})\ge F(R_{\bar u}).
 \label{eq:local-control-projection}
\end{equation}
Here $F$ uses a Gibbs state of the uniform leading $h$, or its uniform
trace. Thus the minimum over all real $u_i$ equals the uniform minimum.
Both the finite-temperature obstruction and the uniform-trace window
bound survive this larger control space. Appendix~\ref{app:local-controls}
gives the proof and microscopic tangent checks. The projection is an
operator identity used in the proof, not an averaging of plotted spectra.

\subsection{Exact activation from a dark four-cell ground state}

The four-cell ring supplies a solvable example of why the distinction
matters. Write $\bm A=\bm T_0+\bm T_2$ and
$\bm B=\bm T_1+\bm T_3$, so that $h=\bm A\cdot\bm B$.
Its ground state has $(j_A,j_B,J)=(3,3,0)$ and energy $-12$;
the first excited multiplet is $(3,3,1)$ at $-11$.
The zero-temperature optimum is $\lambda_0=-2a-3b$, with
$\lambda=-\kappa A'$. The resulting source
$b(Y+3\widetilde D)/\sqrt{20}$ annihilates the ground state exactly,
although it does not commute with $h$.

The first thermally occupied triplet opens two groups of transitions.
Their final sectors are $(3,3,2)$ at $\omega/J_{\rm eff}=2$ and
$(3,1,2),(1,3,2)$ at $\omega/J_{\rm eff}=7$.
Summed over the three initial magnetic components, the squared matrix
elements divided by $b^2$ are respectively $27/50$, $6/25$, and $6/25$.
Consequently,
\begin{equation}
 C_{\min}(\tau)=\frac{51}{50}b^2e^{-1/\tau}
                  +o(e^{-1/\tau}).
 \label{eq:thermal-activation}
\end{equation}
The two groups have leading intensity ratio $9:8$. Reoptimizing the
exchange at each temperature changes this result only beyond the leading
activation order. Exact rational matrix elements reproduce the independently
assembled full-space thermal calculation to relative error
$3\times10^{-15}$ in its leading coefficient. The activation gap, line
positions, and intensity ratio are four-cell results, not universal
long-chain spectral predictions.

\begin{figure*}[t]
\centering
\includegraphics[width=\textwidth]{thermal_evidence.pdf}
\caption{Thermal cancellation and its momentum-acceptance limitation.
(a) Minimum leading positive-frequency coefficient in four- and six-cell
effective periodic chains. The doublet zero is the operator identity,
not a numerical intensity cutoff. (b) Full microscopic two-cell weights
at $T/J_{\rm eff}=1.6$, normalized separately to their values at
$g=\beta/\Dloc=0.005$. Quartet and doublet slopes are $2.02549$ and
$4.00189$; this normalization does not equate their absolute intensities.
(c) Weights averaged over the uniform interval $q_0\pm\delta q$, relative
to their ideal central-momentum values, at $g=0.005$ and
$T/J_{\rm eff}=1.6$. The quartet and doublet central momenta are
$\pi$ and $3\pi/5$, respectively. Exchange parameters remain on their
central-momentum retuning trajectories. Both initial and final microscopic
states lie below $E_0+\Dloc/2$, while the partition function includes
the full spectrum. Lines connect computed points; no spectral or image
smoothing is applied. These are finite-chain results, not an instrument
resolution estimate.}
\label{fig:thermal-evidence}
\end{figure*}

\subsection{Microscopic check and finite momentum acceptance}

Full thermal traces on four- and six-cell effective chains give the
contrast in Fig.~\ref{fig:thermal-evidence}(a). At $\tau=1.6$, the quartet
minima are $0.00991078$ and $0.00968212$, respectively. For six cells the
latter is about 22 times its zero-temperature value, even after
temperature-dependent retuning. This finite-size comparison supports the
operator result without replacing it by a size extrapolation.

To test the perturbative prediction in the original model, we diagonalized
all 1024 states of two coupled pentamers. Two distinct microscopic closing
links require two corresponding effective exchange bonds.
The quartet root is Eq.~\eqref{eq:rank-exact-root}; the doublet uses
$q=3\pi/5$, $\alpha_*=1.09755143$, so both have $Q=\pi$.
At $g=0.005,0.01,0.02$ and $\tau=0.4,1.6$, the retuned quartet
weights have logarithmic slopes $2.07089,2.02549$, compared with
$4.00185,4.00189$ for the doublet. At the weakest coupling the quartet
coefficients differ from their predictions by $1.75\%$ and $1.18\%$.
At $\tau=0.4,g=0.02$ the error reaches $12.25\%$; quantitative
agreement is not claimed throughout the tested interval. These two-cell
calculations test the finite-temperature expansion and order contrast.

We tested the four-cell activation directly in the twenty-spin periodic
microscopic model at $g=0.000625,0.00125,0.0025$. All 256 retained states
were reconstructed by total-spin raising from the 44 low-band states in
the $S^z=0$ sector. The exchange was kept on the same leading
zero-temperature retuning trajectory, $\kappa_0=1.28029750$, rather than
reoptimized for each temperature. At the weakest coupling the first
thermal gap is $1.000318J_{\rm eff}$; the two line-group centroids are
$2.000592J_{\rm eff}$ and $7.000309J_{\rm eff}$. Their unweighted
transition strengths, summed over the initial triplet and divided by
$\beta^2b^2$, are $0.536037$ and $0.482098$, within $0.74\%$ of
$27/50$ and $12/25$. Both approach their predictions as $g$ decreases
[Fig.~\ref{fig:four-cell-activation}(b)].

The full thermal response also exposes two limits of the leading result.
At $\tau=0.1$ the weakest-coupling weight exceeds the prediction by
$25.9\%$: the ground-state residual accounts for $21.0\%$ of the
microscopic total and scales approximately as $\beta^{4.0015}$ over the
three tested couplings. It is retained in
Fig.~\ref{fig:four-cell-activation}(a), not subtracted to reveal an
activation law. At $\tau=1.6$, quasielastic lines originating from split
effective degeneracies contribute $1.68\%$ of the total positive-frequency
weight. A resolved window $\eta=0.05$ separates these lines from the
$O(J_{\rm eff})$ response. Its weakest-coupling coefficients agree with
the fixed-retuning effective predictions to $0.161\%$, $0.229\%$, and
$0.042\%$ at $\tau=0.2,0.4,1.6$; changing $\eta$ from $0.02$ to
$0.1$ leaves the computed window weights unchanged. The window was
introduced as a diagnostic after the threshold sensitivity was detected;
the original unwindowed results are retained. These results validate a
finite-frequency activation mechanism with a finite-coupling crossover,
not an activation law uniformly valid down to $T=0$.

\begin{figure*}[t]
\centering
\includegraphics[width=\textwidth]{four_cell_activation.pdf}
\caption{Direct four-cell microscopic check and low-temperature background.
(a) Resolved thermal weight for twenty physical spins with periodic
boundaries, $q=\pi$, fixed zero-temperature retuning, and
$\omega>0.05J_{\rm eff}$. Both initial and final states are in the
complete 256-state low band. The finite-coupling curves retain the
ground-state background. The effective curve is the full leading thermal
trace at the same retuning; the activation asymptote is
$(51/50)b^2\exp(-J_{\rm eff}/T)$.
(b) Strengths from the first thermally occupied triplet to the two predicted
final-state groups, divided by $\beta^2b^2(27/50)$ and
$\beta^2b^2(12/25)$, respectively. These strengths are summed over initial
magnetic components before thermal weighting. Lines connect calculated
points without smoothing. The omitted-band population has a rigorous
bound, and the frequency-window diagnostic is separate from the original
unwindowed calculation.}
\label{fig:four-cell-activation}
\end{figure*}

The ideal momentum condition is restrictive. We also averaged the
microscopic transition probabilities over a uniform window around each
central momentum, keeping the central retuning fixed. At
$g=0.005,\tau=1.6$, a half-width $0.001\pi$ raises the quartet weight
by $0.094\%$, but the doublet weight by a factor of approximately
$2.93\times10^6$ [Fig.~\ref{fig:thermal-evidence}(c)].
The local doublet form factor has $\partial_qA=1.69998$ at its root,
whereas the quartet root at $q=\pi$ is stationary in momentum.
A fixed window therefore adds an ordinary form-factor background that
can overwhelm the doublet's very small fourth-order central signal.
The observed doublet coupling slope falls to approximately $0.00475$
at that window width. The operator-level distinction at ideal momentum
remains valid, but by itself is not a resolution-independent experimental
diagnostic.

Optimizing the exchange for the entire window reduces, but does not
eliminate, this background. At the stationary quartet root, write
$x=q-\pi$, $A_{qq}=\partial_q^2 A(\alpha_*,\pi)$, and let $B$ be the
positive-frequency thermal weight of $D$. If $\kappa_0$ minimizes the
central response, the remaining source is thermally orthogonal to $D$.
The leading window-dependent dipole amplitude is
$\beta A'(\kappa-\kappa_0)+A_{qq}x^2/2$.
Retuning removes its mean, leaving the variance of $x^2$.
For a uniform half-width $\delta q=\rho\sqrt{\beta}$, at fixed finite
size and scaled temperature, this gives
\begin{align}
 \kappa_{\rm win}&=\kappa_0-\frac{A_{qq}\rho^2}{6A'}+o(1),\notag\\
 \frac{\overline W^{\rm opt}_{+,\eta}}{\beta^2}
 &\longrightarrow C_{\min}+\frac{A_{qq}^2B}{45}\rho^4.
 \label{eq:window-retuning}
\end{align}
Here the regular-frequency window is the one defined above. Keeping
$\kappa_0$ fixed instead gives $1/20$ in place of $1/45$ in the
background coefficient. Thus window retuning removes its mean shift,
not the variation of the form factor across the measured interval.
In the ten-spin ring at $\tau=1.6$, the optimized microscopic
coefficients for $\rho=0.25,0.5,1$ agree with this limit within
$0.30\%$ at $g=0.00125$; deviations increase with coupling.

The chosen doublet root has nonzero $A_q=\partial_q A$, so its leading
window background is instead $A_q^2B(\delta q)^2/3$. Keeping it at
the fourth-order scale requires $\delta q=O(\beta^2)$, whereas the
stationary quartet root permits $O(\sqrt{\beta})$ at its quadratic
scale. These different tolerances reflect the orders of the selected
momentum zeros, not a universal resolution advantage of higher spin.
Furthermore, optimizing beyond the first-order trajectory reduces the
two-cell doublet's central low-band response to numerical zero: only
one singlet--triplet reduced transition amplitude is involved there.
For $\delta q=\rho\beta^2$ its remaining window weight approaches
$\beta^4 A_q^2B\rho^2/3$, with a $0.055\%$ deviation at the weakest
tested coupling. Thus a measured fourth power alone need not identify
an intrinsic central-momentum residual. These shrinking-window tests
do not establish detectability at fixed instrumental resolution or
complete second-order cancellation on longer chains.

\section{Ground-state and boundary controls}

The thermal distinction is supported by separate zero-temperature checks,
detailed in Appendix~\ref{app:ground-controls}. On the same fifteen-spin
periodic microscopic geometry, the retuned quartet and doublet weights
have weak-coupling slopes $1.9869$ and $4.0167$. An even six-cell effective
ring at the same local roots retains the contrast, and converged
twelve-cell periodic MPS results exclude an open-end origin of the
quartet residual at the tested roots. These controls are not a
thermodynamic extrapolation. In particular, the four-cell dark ground
state is a real exception, not a numerical error to be removed.

Open chains have an additional one-spin endpoint contribution, so their
residual cannot be identified entirely with the composite vector.
Microscopic twenty-spin open-chain weights and individual lines approach
the effective predictions at weak coupling, but depart strongly at larger
coupling. The previously calculated $L=40$ full spectra are therefore
retained only as finite-coupling context in
Fig.~\ref{fig:rank-original}; they are not quantitative evidence for the
present retuned weak-coupling mechanism.

\section{Discussion and conclusion}

The thermal response separates operator cancellation from a ground-state
interference zero. Under the stated periodic symmetries, the doublet's
first-order probe can be canceled throughout the retained space; the
quartet's different bond-spin corrections prevent this. Its four-cell
ground state is exactly dark but supports an activated thermal response,
verified in the complete twenty-spin low band.
The exact second moment quantifies the obstruction for the prescribed
dipole--composite probe at general spin. Locality strengthens it to a
size-independent finite-frequency weight bound in the effective uniform
trace. Independent infinitesimal SU(2)-invariant controls inside each cell
cannot improve on uniform retuning: additional operator momenta add
nonnegative weight. Thus the residual is not an artifact of allowing only
one exchange knob.

This interpretation builds on established projection, tensor, and
effective-observable methods. In particular, higher-order scattering operators
are not introduced here for the first time~\cite{waldmann2005manyspin}, nor is
the correction of a cluster interference factor~\cite{mueller2003structure}.
The distinction between symmetry-fixed momentum dependence and
interference involving independent dynamical matrix elements is already
explicit in Ref.~\onlinecite{tabrizi2021selection}, Eqs.~(13)--(16).
Our result determines what remains after optimizing the physically
accessible coefficients of such amplitudes, rather than inferring a
residual simply from the presence of more than one operator.
Conditional excitation processes and thermal redistribution of spectral
weight are also established effects~\cite{exius2010vertex,honecker2016thermal}.
Even the thermal lifting of an accidental scattering zero has an
experimental precedent in magnetovibrational interference~\cite{raymond2018interference}.
The distinction examined here is whether a physical exchange retuning can
remove the corrected response simultaneously across the thermally occupied
low-energy manifold. The surviving vector acts on neighboring effective
spins; it is not a direct rank-two probe or a new excitation species.
The local representation constrains this operator, while ground-state
accessibility determines whether its independence is already visible at
zero temperature.

Several qualifications determine the scope of that distinction. The local
spin permits a tensor but does not ensure that it acts independently on the
many-body ground state, as the four-cell periodic exception demonstrates.
Open endpoints create a separate one-spin residual. The microscopic
same-geometry test is an odd frustrated ring; the contrast survives an even
ring at the same roots, but that check uses the leading effective model.
Converged twelve-cell periodic data rule out a purely endpoint explanation
at the tested quartet root. The finite-temperature commutator proof removes
a dependence on particular even chain lengths, but not the fixed-size
weak-coupling assumption. The weight bound in
Eq.~\eqref{eq:quartet-window-bound} applies to the effective uniform trace,
with that limit taken before the size limit. It is not a thermodynamic
lower bound at arbitrary fixed scaled temperature and does not establish
uniform control of the microscopic expansion with size. The general-spin
results also presuppose a specified nonzero composite probe coefficient.
The local-control extension does not cover anisotropy, changes of local
irrep, or modifications of the leading intercell bond ratios.
In the four-cell microscopic check, the fourth-order
ground-state background eventually dominates the activated response at
fixed nonzero coupling. Higher-order level splittings also generate
quasielastic weight absent from the leading finite-frequency coefficient.
The order of the weak-coupling, low-temperature, and frequency-window
limits therefore matters. Furthermore,
finite momentum acceptance can overwhelm the ideal doublet signal with
ordinary form-factor weight. Window-optimized exchange removes the mean
form-factor shift but leaves its variance; the required shrinking-window
scales depend on the order of the selected momentum zero.
There is no material-specific demonstration of detectability.
The doublet's observed fourth-order behavior along the first-order
trajectory is not an unavoidable floor. Further retuning numerically
cancels the central response in the two-cell example, whose retained
space contains only one singlet and one triplet. A complete second-order
probe and second-order retuning for general chains have not been analyzed.

Within these limits, the pentamer chain gives a quantitative distinction
between an operator-level cancellation and a state-dependent dark point.
Its value is in identifying which inference can be drawn from weak
scattering: a dark zero-temperature spectrum does not by itself imply a
thermal selection rule or absence of the corresponding excitations.
Using the distinction experimentally requires control of momentum
acceptance and nonleading backgrounds in addition to the effective
Hamiltonian and local spin representation.

\begin{acknowledgments}
[Author, affiliation, funding, and acknowledgment details are retained for
completion by the authors.]
\end{acknowledgments}
\section*{Data and code availability}
The working project contains the calculation scripts, raw numerical tables,
protocol records, and SHA-256 manifests used here. It has not been deposited in
a public repository, and no accession identifier is claimed.
\appendix
\section{Exact second moment for the dipole--composite probe}
\label{app:general-moment}

This appendix assumes a uniform spin-$s$ nearest-neighbor periodic chain
with even $N\ge4$, $s\ge1/2$, and the probe form specified before
Eq.~\eqref{eq:general-moment}. Set $c=s(s+1)$ and $d=2s+1$.
On one bond, the quadrupole--spin vector obeys
\begin{equation}
 X_{ij}=\tfrac12\{h_{ij},D_{ij}\}+\tfrac{c}{3}D_{ij},
 \qquad D_{ij}=T_i^z-T_j^z.
\end{equation}
In the $J-1\to J$ channel its coefficient relative to $D_{ij}$ is
$r_J=J^2/2-2c/3$. The squared $D^z$ matrix elements summed over magnetic
components are $J(d^2-J^2)/3$, and the corresponding energy difference is
$J$. Including both transition directions, define
$w_J=2J^3(d^2-J^2)/3$ for $1\le J\le d-1$. Thus the Gram entries of
$C_D=[h_{ij},D_{ij}]$ and $C_X=[h_{ij},X_{ij}]$ are sums of
$w_J$, $w_Jr_J$, and $w_Jr_J^2$:
\begin{align}
 G_{DD}&=\frac{d^2(d^2-1)^2}{18},\notag\\
 \frac{G_{DX}}{G_{DD}}&=\frac{d^2-4}{12},\qquad
 \frac{G_{XX}}{G_{DD}}=\frac{7(d^2-4)^2}{360}.
 \label{eq:general-bond-gram}
\end{align}
These finite sums follow from power-sum identities. We verified the
identities coefficient by coefficient through their finite differences,
rather than fitting values at selected $s$. The minimized bond norm is
$d^2(d^2-1)^2(d^2-4)^2/1440$.
For $s=3/2$, the Gram matrix is
\begin{equation}
 \begin{pmatrix}G_{DD}&G_{DX}\\G_{DX}&G_{XX}\end{pmatrix}
 =\begin{pmatrix}200&200\\200&560\end{pmatrix},
 \label{eq:thermal-gram}
\end{equation}
with minimized bond norm $360$.

The adjacent-bond contribution on sites $0,1,2$ is
\begin{equation}
 Z_{012}=[h_{12},X_{01}]-[h_{01},X_{12}].
\end{equation}
With repeated Cartesian indices summed, its explicit decomposition is
\begin{align}
 Z_{012}/i={}&\epsilon_{abc}
  (Q_0^{zb}T_1^cT_2^a+T_0^aT_1^cQ_2^{zb})\notag\\
 &-T_0^a(\epsilon_{bzc}Q_1^{ca}
                  +\epsilon_{azc}Q_1^{cb})T_2^b.
 \label{eq:three-site-commutator}
\end{align}
The three terms have local tensor-rank patterns $QTT$, $TTQ$, and
$TQT$, respectively. They are mutually orthogonal because
$\operatorname{tr}(T^aQ^{bc})=0$. Rotational invariance and the traces of
$(T^z)^2$ and $(T^z)^4$ give
\begin{align}
 \operatorname{tr}(T^aT^b)&=t\delta_{ab},\quad t=dc/3,\notag\\
 \operatorname{tr}(Q^{ab}Q^{ef})&=v(\delta_{ae}\delta_{bf}
 +\delta_{af}\delta_{be}-\tfrac23\delta_{ab}\delta_{ef}),\notag\\
 v&=dc(4c-3)/60.
\end{align}
Contracting the Levi-Civita symbols in
Eq.~\eqref{eq:three-site-commutator} gives $20/3$, $20/3$, and $20$
times $vt^2$ for those three patterns. Therefore
\begin{equation}
 \|Z_{012}\|_{\rm HS}^2=\frac{100}{3}vt^2
                      =\frac{5}{81}d^3c^3(4c-3).
 \label{eq:three-site-norm}
\end{equation}

All one-site partial traces of $C_D,C_X$ and all two-site partial traces
of $Z$ vanish. Distinct two-site and three-site supports are consequently
orthogonal in the full uniform trace. For a periodic chain with $N\ge4$,
there are $N$ distinct supports of each type. The alternating signs give
\begin{equation}
 M_2(u)=\frac{\|bC_X+uC_D\|_{\rm HS}^2}{\ell d^2}
       +\frac{b^2\|Z\|_{\rm HS}^2}{\ell d^3}.
\end{equation}
Completing the square with Eqs.~\eqref{eq:general-bond-gram} and
\eqref{eq:three-site-norm} yields Eq.~\eqref{eq:general-moment}.
The first contribution alone is the earlier lower bound. The second
cannot be tuned away by changing $u$.
Keeping only the two-site contribution already gives the quartet bound
\begin{equation}
 \min_u\frac{\|[h,R]\|_{\rm HS}^2}{4^N}
 \ge\frac{360b^2}{5\times16}=\frac92b^2,
 \label{eq:thermal-bound}
\end{equation}
which the exact result sharpens to $(197/16)b^2$.

The local identities and partial traces were additionally checked with
exact rational matrices for $s=1/2,1,3/2,2,5/2,3$. Separate orthonormal
Cartesian-spin Kronecker matrices on the full $N=4,6$ chains, without
importing the rational proof routines, verify the whole quadratic for
$s=1/2,1,3/2$ to floating-point accuracy. These checks support the
algebraic derivation; they do not extend its assumptions to other probe
structures, open chains, odd rings, or nonuniform exchange.

\subsection{A size-uniform positive-frequency weight bound}
\label{app:weight-bound}

Use the normalized Hilbert--Schmidt norm on each operator's support,
$\|O\|_2^2=d^{-|\operatorname{supp}O|}\operatorname{Tr}O^\dagger O$.
The full uniform-trace spectral measure of $R$ is symmetric in frequency.
Define its moments $M_{2k}$ including both frequency signs; $M_0$ also
includes equal-energy matrix elements. Local tensor orthogonality gives
\begin{equation}
 M_0=\frac{1}{\ell}\left[\frac{c}{3}u^2
          +\frac{c^2(4c-3)}{27}b^2\right].
 \label{eq:uniform-M0}
\end{equation}
For $s\ge1$, comparison with Eq.~\eqref{eq:general-moment} yields
$M_2\ge(3c/2)M_0$. To see positivity without optimizing an irrational
eigenvalue, rescale $b,u$ by the square roots of their coefficients in
$\ell M_0$. The matrix for $\ell[M_2-(3c/2)M_0]$ in these coordinates has
diagonal entries $(61c-42)/30$ and $7c/6$, and determinant
$c(107c-54)/180$. All are positive for $c\ge2$.

For the fourth moment write
\begin{align}
 [h,R]&=\frac{1}{\sqrt{\ell N}}\sum_i(-1)^i B_i,\notag\\
 B_i&=bC_{X,i}+uC_{D,i}+bZ_{i,i+1,i+2}.
\end{align}
The two-site and three-site terms are orthogonal as above, giving
$\|B_i\|_2^2=\ell M_2$. Only the four bonds meeting sites $i,i+1,i+2$
contribute to $G_i=[h,B_i]$. Since $\|h_{ij}\|_{\rm op}=c$,
$\|G_i\|_2\le8c\|B_i\|_2$. Each traceless $G_i$ is supported within
the five-site interval from $i-1$ to $i+3$. Two disjoint intervals have
zero uniform-trace inner product. At most nine translated intervals,
including itself, overlap a given interval; for $N=4,6,8$ the number is
even smaller. Bounding the remaining inner products by their norms gives
\begin{equation}
 M_4\le\frac{9}{\ell}\|G_i\|_2^2
      \le576c^2M_2.
 \label{eq:uniform-M4-bound}
\end{equation}
This is a locality bound, not a fit to numerical moments.

Separate the spectral measure into $|\omega|\le\nu$ and
$|\omega|>\nu$. The low-frequency contribution to $M_2$ is at most
$\nu^2M_0$. Cauchy--Schwarz for the remaining measure therefore gives
\begin{equation}
 C^{\rm tr}_{>\nu}\ge
 \frac{[M_2-\nu^2M_0]_+^2}{2M_4},
 \label{eq:window-moment-inequality}
\end{equation}
where $[x]_+=\max(x,0)$ and the factor two accounts for the two frequency
signs. Taking $\nu=\sqrt{3c/4}$ and using the preceding inequalities,
\begin{align}
 C^{\rm tr}_{>\sqrt{3c/4}}
 &\ge\frac{M_2}{4608c^2}\notag\\
 &\ge\frac{b^2(4c-3)(86c-27)}{3732480\ell}>0
 \quad(s\ge1,\ b\ne0).
 \label{eq:general-window-bound}
\end{align}
This holds uniformly over every real retuning $u$ and every finite even
$N\ge4$. It is the ordered uniform-trace limit of the effective model,
not an assertion about arbitrary finite scaled temperature or the full
microscopic infinite-temperature spectrum.

Full-chain matrices at $s=1,3/2$ and $N=4,6$ verify the zeroth-, second-,
and fourth-moment inequalities. Complete four-site diagonalization also
checks Eq.~\eqref{eq:window-moment-inequality} from actual transition
weights. Minimizing the quartet window weight gives $0.350008b^2$,
about 1842 times its conservative lower bound. Thus the lower bound is
useful as a size-uniform nonvanishing statement, not as a quantitative
estimate of the scattering intensity.

\section{Extension to independent cell-local controls}
\label{app:local-controls}

We first connect a microscopic control to the enlarged probe space in
Eq.~\eqref{eq:local-control-probe}. Add
$\delta H_{\rm loc}=\beta\sum_i W_i$, where each bounded Hermitian $W_i$
is an SU(2) scalar confined to cell $i$. The reference ground multiplet is
an isolated single spin-$s$ irrep. Schur's lemma gives
$P_iW_iP_i=w_iP_i$, a constant rather than an effective spin interaction.
The first-order derivative of the local vector probe is again a vector.
Because rank one occurs once in the operator space of a single spin-$s$
irrep, its projection must be a multiple of $T_i^z$. At $q=\pi$ the
probe is Hermitian, so this multiple is real. Changes in the projected
end-spin coefficients enter $\beta V$ only at order $\beta^2$;
likewise, the change of $b$ affects the physical first-order source only
at order $\beta^2$. The leading Hamiltonian is consequently still the
same uniform $h$, and all such controls are contained in the real
one-spin space of Eq.~\eqref{eq:local-control-probe}. Cell reflection
symmetry need not be preserved by $W_i$. This is a tangent-space result
about the reference root, not a finite displacement along a family of
different microscopic models.

Let $U$ translate by one effective site. For eigenstates of $h$, define
the positive semidefinite sesquilinear form
\begin{equation}
 (A,B)_{\tau,\nu}=\sum_{\epsilon_f-\epsilon_i>\nu}
 p_i A_{fi}^*B_{fi},\qquad \nu\ge0,
 \label{eq:window-inner-product}
\end{equation}
where $p_i$ are Gibbs weights or $1/d^N$. All degeneracies are included
as complete blocks. Since $[U,h]=0$, conjugation by $U$ is an isometry
of this form. Distinct operator momenta are orthogonal even when the
form has a null space. The $\pi$-momentum projection is
\begin{equation}
 \Pi_\pi A=\frac1N\sum_{m=0}^{N-1}(-1)^mU^m A U^{-m}.
\end{equation}
For even $N$, it obeys
$\Pi_\pi R_{\{u_i\}}=R_{\bar u}$. Pythagoras in
Eq.~\eqref{eq:window-inner-product} gives
Eq.~\eqref{eq:local-control-projection}, and hence
\begin{equation}
 \min_{\{u_i\}\in\mathbb R^N}F(R_{\{u_i\}})
 =\min_{u\in\mathbb R}F(R_u).
 \label{eq:local-control-optimum}
\end{equation}
The equality uses the availability of uniform controls in the enlarged
space. Any more restricted microscopic control set still satisfies the
same lower bound. The nonzero derivative $A'$ makes the uniform direction
accessible in the pentamer example. This argument applies at every
positive scaled temperature and to the uniform-trace window bound;
it does not require a nondegenerate spectrum.

The second moment provides an additional explicit check. On bond
$(i,i+1)$ the dipole commutator coefficient becomes
$(u_i+u_{i+1})/2$. Repeating the support-orthogonality calculation of
Appendix~\ref{app:general-moment} gives
\begin{align}
 M_2(\{u_i\})={}&M_{2,\min}+\frac{(d^2-1)^2}{18\ell N}\notag\\
 &\times\sum_i\left[\frac{u_i+u_{i+1}}2+
              \frac{b(d^2-4)}{12}\right]^2.
 \label{eq:nonuniform-moment}
\end{align}
An alternating component of $u_i$ adds a conserved total $T^z$ to the
probe and leaves this moment unchanged. For the same reason, the earlier
$M_2/M_0$ inequality cannot be applied unchanged to an arbitrary
nonuniform control. The finite-frequency lower bound follows instead
from the positive operator-momentum decomposition above.

For all ten pairwise exchange directions inside the isolated pentamer,
the derivative $-P(W_iGB+BGW_i)P$, with
$G=Q(H_{\rm cell}-E_0)^{-1}Q$, has a non-dipole Frobenius norm below
$2.4\times10^{-15}$. Central differences at step $5\times10^{-5}$
agree within $2.8\times10^{-10}$. Full four-site spectra for $s=1,3/2$
at $\tau=0.2,1.6$ and uniform trace verify
Eq.~\eqref{eq:local-control-optimum} for $\nu=0$ and
$\sqrt{3s(s+1)/4}$, with maximum difference $1.7\times10^{-16}$ at
$b=1$. Full-chain matrices at $N=4,6$ verify
Eq.~\eqref{eq:nonuniform-moment} within $1.5\times10^{-14}$.

\section{Microscopic tests and finite-chain controls}
\label{app:ground-controls}

\subsection{Different powers on the same periodic geometry}

The most direct microscopic comparison uses the same $N=3$, $L=15$ periodic
geometry for both local spins. Table~\ref{tab:rank-roots} specifies simple
roots at different physical momenta but the same effective momentum
$Q=2\pi/3$. The doublet root extends the positive-$\alpha$ interval beyond
the previous $\alpha\le1$ survey; no microscopic interaction is added.
The odd periodic geometry is frustrated, so this comparison tests the
finite-system mechanism rather than an infinite unfrustrated chain.

\begin{table}[t]
\caption{Parameters fixed before the microscopic three-cell comparison.
The local gap is evaluated at the root. $C_0=\|k\|^2$ and $\Cmin$ are
effective predictions for the coefficient of $\beta^2$, not fits to the
microscopic weights.}
\label{tab:rank-roots}
\begin{ruledtabular}
\begin{tabular}{lrr}
 & Quartet & Doublet\\
$s$ & $3/2$ & $1/2$\\
$q/\pi$ & $14/15$ & $8/15$\\
$\alpha_*$ & $-0.38103454$ & $1.45352917$\\
$\Dloc$ & $0.10336638$ & $1.10511693$\\
$\kappa_{\mathrm{opt}}$ & $-9.45913710$ & $-0.17229706$\\
$C_0$ & $0.54511625$ & $0.000474008$\\
$\Cmin$ & $0.14202789$ & $<10^{-30}$\\
\end{tabular}
\end{ruledtabular}
\end{table}

We diagonalized the microscopic $S^z=1/2$ sector of dimension 6435 at
$g=0.005,0.01,0.02,0.04$, both at the fixed local root and with its
predicted linear retuning. There are two degenerate ground states in this
sector; their equal average excludes the entire ground eigenspace from the
inelastic weight. All ten quartet or one doublet low-band excited states
in this sector were included, with additional states above the cutoff used
to check band completeness. The largest eigen-residual was below
$3.1\times10^{-13}$.

The retuned quartet weight remained quadratic at weak coupling
[Fig.~\ref{fig:rank-evidence}(a)]. At the smallest $g$, its coefficient
$\Wlow/\beta^2=0.14091678$ differs from the independent prediction
$0.14202789$ by $0.7823\%$. A log--log fit over the three smallest couplings
gave a slope $1.9869$. For the doublet the corresponding slope was $4.0167$
[Fig.~\ref{fig:rank-evidence}(b)]. Its $\Wlow/\beta^4$ varied from
$1.0694\times10^{-4}$ to $1.1290\times10^{-4}$ across all four couplings.
These finite-interval slopes test perturbative orders, not critical exponents
or statistical scaling laws. The doublet result supports
Eq.~\eqref{eq:rank-doublet}, but does not prove that a fourth-order floor
survives further retuning.

The cancellation was also checked at the operator level. In the full
eight-dimensional retained doublet space, the relative norm of
$K_q+\kappa_{\mathrm{opt}}A'D_Q$ was $4.32\times10^{-15}$.
The analogous quartet operator did not vanish in its full 64-dimensional
space. A restart with a different random vector, larger Krylov space, and
tighter tolerance reproduced the two weakest doublet weights to relative
differences below $1.3\times10^{-9}$.

\begin{figure*}[t]
\centering
\includegraphics[width=0.98\textwidth]{rank_compensation.pdf}
\caption{Retuned-weight comparison and validation limits.
(a,b) Absolute microscopic low-band weights on the same three-cell periodic
geometry, at the parameters in Table~\ref{tab:rank-roots}. Dashed curves keep
$\alpha=\alpha_*$; solid curves use $\alpha_*+\kappa_{\mathrm{opt}}\beta$.
The vertical ranges differ, but no independent intensity normalization is
applied. Points are calculated values; segments only connect them.
(c) Leading quartet minimum $\Cmin$ at the exact $q=\pi$ root for even
effective chains. ED denotes exact diagonalization, OBC/PBC denote open/periodic
boundaries. The accepted DMRG point is an effective $N=20$ open chain;
the $N=40$ cross failed the preset energy-convergence criterion and is not
used as converged evidence. No thermodynamic fit is drawn.
(d) Microscopic $L=20$ open-chain comparison of the retuned leading probe
with the line-weight vector and integrated weight. The errors at larger
$g$ reveal failure of the leading theory. All panels use calculated data
without smoothing; logarithmic axes in (a,b,d) change only the display.}
\label{fig:rank-evidence}
\end{figure*}

\subsection{Even periodic chains and the open-end contribution}

We first kept precisely the two local roots and the effective momentum
$Q=2\pi/3$ of Table~\ref{tab:rank-roots}, but doubled the ring to six cells.
The quartet effective minimum was $0.06641159$, while the doublet minimum
was below $4\times10^{-32}$. The quartet's optimum changed to
$\kappa=-5.56848$, reflecting its dependence on the many-body ground state;
the doublet's operator-level optimum remained $-0.17229706$.
Thus the contrast also occurs without odd-ring frustration in the effective
model. This is not a microscopic thirty-spin calculation.
At the same quartet root, the accepted $N=12$ periodic MPS gives
$\Cmin=0.08315075$ and $\kappa_{\mathrm{opt}}=-5.67393$
[Fig.~\ref{fig:rank-controls}(d)]. Increasing its bond dimension from
384 to 512 changes this weight by only $6.15\times10^{-7}$ relatively.
Its scalar contribution is removed to within $1.4\times10^{-17}$ in
the residual norm squared, again identifying a composite rather than
endpoint source. These size-dependent coefficients are not extrapolated.

At the exact quartet root in Eq.~\eqref{eq:rank-exact-root}, periodic
effective chains with $N=6,8,10$ gave
$\Cmin=0.000443300,0.000649072,0.000751090$, respectively.
However, $N=4$ gave zero to numerical precision
[Fig.~\ref{fig:rank-evidence}(c)]. This exception rules out attributing an
unavoidable nonzero residual to the local spin alone.

Periodic DMRG extends this test to $N=12$, where bond dimension 512 gives
$\Cmin=0.00080811794$. Relative to bond dimension 384, the coefficient
changes by $3.52\times10^{-5}$ and the energy per cell by
$2.97\times10^{-9}$, passing the original convergence thresholds.
The energy variance per cell is $6.73\times10^{-10}$, and the bond-energy
spread is $1.41\times10^{-8}$. These periodic results exclude an endpoint
origin for this finite-chain residual without establishing a thermodynamic
lower bound.
The $N=16$ calculation at bond dimension 768 had a stable weight but an
energy change per cell of $1.304\times10^{-8}$, above the preset
$10^{-8}$ gate. It is marked unaccepted in Fig.~\ref{fig:rank-controls}(c)
and is not used to extend the convergence claim.

To distinguish bulk-link physics from endpoints, we separated $K_q$ into the
one-spin and quadrupole-times-spin pieces of Eq.~\eqref{eq:rank-link} before
projecting each perpendicular to $d$. For all periodic cases, the perpendicular
one-spin norm squared was below $2\times10^{-30}$. The nonzero residuals are
therefore due to the composite source. Open chains do not share this property.
At $N=10$ OBC, the one-spin, composite, and interference contributions to
$\Cmin$ were $0.000286554$, $0.001244067$, and $-0.000513469$.
Their sum is $0.001017151$. Calling this whole open-chain residual a pure
composite response would be incorrect.

Retuning also changed which excitations dominated the remaining effective
weight. For $N=10$ PBC, the lowest-line share decreased from $0.94007$ at the
fixed root to $1.04\times10^{-5}$ at the integrated-weight optimum. The mean
energy of the residual, in units of $\beta$, rose from $0.09456$ to $0.85896$.
These numbers describe an effective finite chain. They do not imply a bulk
continuum, and the integrated optimum is not exactly the value that cancels
the lowest line alone.

An independent MPS calculation reproduced the $N=10$ open-chain ED minimum
to relative error $2.4\times10^{-6}$ at bond dimension 128. Refinement to
bond dimension 256 gave an accepted $N=20$ minimum $0.00100937976$.
The $N=40$ weight appeared stable but its energy change per cell,
$1.0814\times10^{-8}$, remained above the preset $10^{-8}$ threshold.
It is reported as unaccepted, not used to infer a nonzero thermodynamic limit.

\subsection{Microscopic validity at the exact quartet root}

\begin{figure*}[t]
\centering
\includegraphics[width=0.98\textwidth]{rank_controls.pdf}
\caption{Local mechanism and periodic-chain controls.
(a) Simple local roots at 16 common effective momenta, with
$q=(Q+4\pi)/5$ for quartets and $q=(Q+2\pi)/5$ for doublets. Segments connect
calculated roots; this is not a many-body phase diagram.
(b) Exact channel coefficient $r_J$ in Eq.~\eqref{eq:bond-channel}, not a
spectral intensity. Different quartet channels cannot all be compensated by
the same scalar shift at the link-operator level.
(c) Quartet $\Cmin$ at $q=\pi$ on even periodic effective chains.
The $N=12$ DMRG point passed the original convergence tests; the $N=16$
cross did not and is not converged evidence. No bulk fit is drawn.
(d) Effective chains at the same roots and $Q=2\pi/3$ as
Table~\ref{tab:rank-roots}; the quartet optimum is recomputed for each size.
The doublet dashed line is the first-order operator-cancellation result,
with ED dots at $N=3,6$, not a microscopic finite-$\beta$ zero.
Panels (c,d) retain the normalization $1/\sqrt{5N}$ but use different physical
momenta and local roots; they are not one common size series.
All weights are computed without smoothing or reflection averaging.}
\label{fig:rank-controls}
\end{figure*}

Four open pentamers provide a second microscopic test at $q=\pi$, without
the frustration of the three-cell ring. The effective prediction for their
retuned coefficient is $\Cmin=0.00067813935$.
At $g=0.000625$, microscopic diagonalization gave $0.00067904306$, an
integrated-weight error of $0.1333\%$. The relative Euclidean error of the
sorted low-band line-weight vector was $1.2254\%$.
Doubling $g$ to $0.00125$ approximately doubled that line error to
$2.4625\%$; an independent solver restart at the smaller $g$ agreed to
relative line error $3.56\times10^{-9}$.

The agreement is restricted to weak coupling. At $g=0.02$ and $0.04$, the
line-vector errors were $46\%$ and $109\%$, and the integrated-weight errors
were $14\%$ and $50\%$ [Fig.~\ref{fig:rank-evidence}(d)]. These results
prevent extending the leading composite-probe interpretation quantitatively
to the much stronger couplings of the original full $L=40$ spectra.

\begin{figure*}[t]
\centering
\includegraphics[width=0.49\textwidth]{fig6_fullq_positive.pdf}\hfill
\includegraphics[width=0.49\textwidth]{fig6_fullq_negative.pdf}
\caption{Previously calculated complete connected spectra for $L=40$ open
chains, retained without changing the two selected panels.
$J_2/J_1=0.5$ (left), $-0.5$ (right), and $J_3/J_1=0.5$.
The triangular time window at $TJ_1=60$ has width $0.09277J_1$;
the production time step and bond dimension are $\Delta tJ_1=0.05$ and 192.
Both use the same display coordinate
$u(s)=\operatorname{asinh}(s/0.02)/\operatorname{asinh}(6/0.02)$ for
$s=J_1S_T^{zz}\in[0,6]$. This monotonic display does not change the spectral
arrays or their resolution. No Gaussian smoothing, momentum interpolation,
or reflection averaging is applied. The 41 momenta were independently
calculated. These finite-coupling maps are contextual, not quantitative
evidence for the weak-coupling retuned residual.}
\label{fig:rank-original}
\end{figure*}

\section{Local root and numerical protocols}

For the quartet root at $q=\pi$, the highest-weight $M=3/2$ state has
single-down-spin amplitudes $(u,v,w,v,u)$, with $w=-2(u+v)$ from the
highest-weight condition. The zero of the projected alternating probe gives
$v^2=1/8$ after normalization. The relevant eigenvector ratio is
$x=u/v=(\sqrt7-2)/3$, with
$\alpha=-x(3+2x)/(1-x^2)$, giving Eq.~\eqref{eq:rank-exact-root}.
All local gaps and resolvents were also obtained from direct diagonalization
of the 32-dimensional pentamer Hamiltonian. No reduced transition amplitudes
were fitted to the microscopic spectra.

Microscopic Hamiltonians were built in fixed-$S^z$ bases. The three-cell
comparison used a real symmetric sparse eigensolver with tolerance
$10^{-13}$, at least 50 Krylov vectors, and four extra eigenpairs above the
expected retained band. The numerical acceptance tests required two ground
states in the chosen sector, a complete retained band, and eigen-residuals
below $10^{-10}$. The anticipated small-coupling slope intervals were
$[1.8,2.2]$ and $[3.7,4.3]$; the quartet's weakest-coupling coefficient was
required to agree within $5\%$. These are deterministic validation tests,
not uncertainty intervals. All passed. The parameter and momentum search
used to choose the roots was exploratory; the subsequent microscopic tests
used those fixed predictions.
The restart check used tolerance $10^{-14}$, 120 Krylov vectors, a different
seed, and six extra eigenpairs above the retained count at $g=0.005,0.01$.
It used the same Hamiltonian assembly and is a numerical-sensitivity check,
not an independent computational method.

For the four-cell quartet test, the $S^z=0$ microscopic basis has 184756
states. We obtained 45 eigenpairs, including all 44 retained-sector states
(the ground state and 43 excited states), with tolerance $10^{-13}$.
The two smallest couplings used 140 Krylov vectors; a repeat at
$g=0.000625$ used 180 and a different random starting vector. The line-vector
error compares weights in increasing-energy order, including a zero
ground-state entry, after division by $\beta^2$. It is not a test of a
broadened line shape. The strongest-coupling points were diagnostic tests,
not part of the weak-domain acceptance claim.

Effective-chain ED used the $S^z=0$ sector, whose dimension reaches 116304
for ten spin-$3/2$ sites. The norm and overlap in
Eq.~\eqref{eq:rank-minimum} were evaluated directly on the ground-state
sources, without summing an incomplete set of excited states. The lowest
line and the energy moment were checked separately. The periodic scalar
source and the sum of the open-chain decomposition supplied internal
operator checks.

The longer open effective chains were computed by DMRG~\cite{white1992,schollwoeck2011}
using ITensors/ITensorMPS~\cite{fishman2022}
with conserved $S^z$, explicit spin-$3/2$ local operators, and MPO contractions
of $D^\dagger D$, $K^\dagger K$, and $D^\dagger K$. Thus there is no additional
truncation from constructing an applied source MPS. Refinement used bond
dimensions $128,192,256$, with 24, 16, and 16 sweeps, cutoff $10^{-13}$,
and eigensolver tolerance $10^{-13}$. Larger-bond runs started from the
previous state without decreasing its bond dimension. Acceptance required
relative change of $\Cmin$ below $1\%$, energy change per cell below
$10^{-8}$, energy variance per cell below $10^{-8}$, and maximum local
$|\langle T_n^z\rangle|<10^{-4}$. The final $N=20$ variance was
$1.21\times10^{-9}$ and its maximum local moment was $1.17\times10^{-9}$.
The earlier unconverged runs and the failed $N=40$ refinement are retained
with their original criteria.

Periodic DMRG used the same operators with the closing bond included in the
MPO. The $N=6$ source norm was cross-checked against ED. Bond dimensions
$128,256,384$ were followed by 512 for $N=12$, with cutoff $10^{-13}$ and
16 warm-start sweeps at each refinement. In addition to the preceding
energy and weight criteria, acceptance required
$|\langle\bm T_{\mathrm{tot}}^2\rangle|<10^{-4}$, bond-energy spread below
$10^{-5}$, and scalar-removal and energy-reconstruction identities within
$10^{-10}$. The first $N=12$ batch failed the energy-change test at 384;
the subsequent 512 calculation passed without changing that threshold.
All sizes here count effective spin-$3/2$ sites, not microscopic spins.
The bounded refinement also tested $N=16$ at dimensions $256,512,768$;
the last run failed the unchanged energy-change threshold and no larger
periodic calculation was started.

For the $Q=2\pi/3$ quartet check, the same periodic ground states can be
reused: changing the local root multiplies the leading uniform Heisenberg
Hamiltonian by the positive scalar $c_1(\alpha)^2$ without changing its
eigenvectors. The probe coefficients were recomputed at the new root,
including their complex Fourier phases. Connected norms use
$\langle D^\dagger K\rangle-\langle D\rangle^*\langle K\rangle$,
and the exchange retuning remains real. The six-cell MPO result reproduced
ED to relative error $4.5\times10^{-15}$. The twelve-cell weight passed the
$1\%$ bond-dimension-change test, and the imaginary part of the normalized
source overlap was below $10^{-8}$. This is an effective-model measurement,
not microscopic sixty-spin diagonalization.

\section{Thermal traces and exact local checks}
\label{app:exact-local-check}

The four- and six-cell thermal effective calculations included every
$S^z$ sector, with total dimensions $4^N$ or $2^N$. For six quartet
cells the largest block had dimension 580. Real and imaginary operator
components were transformed into the complete energy bases, and all
transitions with $|\epsilon_f-\epsilon_i|<10^{-9}$ were treated as
equal-energy transitions. Checks included Hermiticity, trace sum rules,
eigen-residuals below $10^{-11}$, detailed balance, and agreement between
direct residual norms and the thermal least-squares expression. A separate
256-dimensional tensor-product assembly reproduced the four-cell result.
Random unitary changes within each degenerate eigenspace left the thermal
weights invariant. Both floating-point assemblies used the same dense
diagonalization library; they are not independent solution methods.

\paragraph*{Exact microscopic prefactors.}
At the quartet root, an unnormalized highest-weight cell state in the
one-down-spin basis, ordered by the position of that spin, is
\begin{equation}
 v=(1,\ 2+\sqrt7,\ -6-2\sqrt7,\ 2+\sqrt7,\ 1)^{\mathsf T}.
\end{equation}
Its squared norm is $88+32\sqrt7$ and its energy is
$E_5=-(5+4\sqrt7)/12$. Total-spin lowering constructs the quartet
projector $P_{\rm cell}$. We evaluated the reduced resolvent
$G=(H_{\rm cell}-E_5+P_{\rm cell})^{-1}-P_{\rm cell}$ separately in all
six magnetization blocks, using exact arithmetic in $\mathbb Q(\sqrt7)$.
An exact $LDL^{\mathsf T}$ factorization of
$H_{\rm cell}-E_5-\tfrac18(1-P_{\rm cell})+P_{\rm cell}$
has strictly positive diagonal pivots in every block. This certifies
that the quartet is the ground multiplet and that $\Dloc>1/8$.
The projected tensor in Eq.~\eqref{eq:rank-local-tensor} is
$s_0\delta_{z\mu}I+t_0Q^{z\mu}$, with
\begin{align}
 s_0&=\frac{161590\sqrt7-361001}{180792},\notag\\
 t_0&=\frac{169277-70390\sqrt7}{45198}.
\end{align}
Both endpoint choices and all $z,+,-$ components satisfy this identity
exactly. Multiplication by $-c_1$ gives $a,b$ in
Eq.~\eqref{eq:exact-local-coefficients}. The elementary bounds
$2.645<\sqrt7<2.646$ certify $b>20086103/162712800>0$.
Comparison with the separately assembled floating-point two-cell
resolvent gives an absolute coefficient difference below
$3.0\times10^{-16}$. This is a check of the same perturbative formula,
not a distinct many-body approximation.

The other exact checks used rational arithmetic in the unnormalized local basis
$|k\rangle=(T^-)^k|s,s\rangle$, where
$T^+|k\rangle=k(2s-k+1)|k-1\rangle$ and
$\langle k|k\rangle=k!(2s)!/(2s-k)!$.
This gives exact zero tests for the partial traces in
Eq.~\eqref{eq:thermal-bound}, with no numerical threshold.
For the four-cell activation calculation, the two sublattice spins were
coupled from $j_A=j_B=3$. In an unnormalized product basis, a highest-weight
total-spin-$J$ state has coefficients $c_k$ at sublattice lowering degrees
$(k,6-J-k)$, obeying
\begin{equation}
 \frac{c_{k+1}}{c_k}
 =-\frac{(6-J-k)(J+1+k)}{(k+1)(6-k)}.
\end{equation}
This recurrence, followed by total-spin lowering, constructs the ground
singlet and first triplet without a floating-point eigensolver.
Acting with $Y+3\widetilde D$ annihilates the singlet. On the triplet,
the residual separates into final eigenvalues $-9$ and $-4$ of $h$;
projection by $\bm A^2$ resolves the two $-4$ channels. Their exact norms
give Eq.~\eqref{eq:thermal-activation}.

The ten-spin microscopic thermal calculation used all 11 $S^z$ blocks,
whose largest dimension is 252. The full partition function was retained.
Both initial and final states in $W_+$ were restricted to
$E-E_0<\Dloc/2$; there were 16 retained quartet or 4 retained doublet
states, as required. Eigen-residuals were below $3.4\times10^{-15}$,
and omitted-band thermal populations were below $1.1\times10^{-57}$
for the tested parameters. The temperature was
$T=\tau\beta c_1(\alpha_*)^2$, including on the retuned trajectory.

For the finite momentum interval we formed the low-band matrix elements
of each microscopic $S_j^z$. At each quadrature point the complex Fourier
amplitudes were combined before taking their absolute squares; this
avoids subtracting nearly equal contributions to a tiny central intensity.
Uniform-window integrations with 16 and 32 Gauss--Legendre nodes agreed
within $10^{-7}$ relatively. Zero-width values reproduced the separate
microscopic calculation within $1.2\times10^{-10}$. This is integration
over a specified momentum acceptance, not Gaussian broadening or an
image-processing operation.

For twenty microscopic spins, the sparse $S^z=0$ Hamiltonian has dimension
184756. We computed 48 lowest eigenpairs and verified that precisely 44
lie below $E_0+\Dloc/2$. Within this low band, diagonalization of
$(S^+V)^\dagger(S^+V)$ resolves total spin. Rediagonalizing the Hamiltonian
within each spin sector and applying normalized total-spin raising gives
all positive-$S^z$ states; spin reversal supplies the negative sectors.
The multiplicities of $J=0,\ldots,6$ are $4,9,11,10,6,3,1$,
giving 256 states. Every raised sector was checked against an independently
assembled microscopic Hamiltonian. Residuals were below $10^{-12}$ and
orthogonality errors below $10^{-11}$. The same procedure at ten spins
reproduced the full-spectrum thermal weight to relative error below
$10^{-7}$ before the twenty-spin run.
At the weakest coupling, an independent restart with a different random
vector, a larger Krylov space (220 instead of 180), and tighter tolerance
changed the total, ground-initial, and excited-initial weights by at most
$8.0\times10^{-8}$ relatively across the tested temperatures and both
frequency definitions. This is solver sensitivity testing, not a distinct
many-body method.

The omitted states have a rigorous thermal bound. Relative to uncoupled
cells at $\alpha_*$, the perturbation norm is at most
$v=(3/4)[N|\beta|+2N|\alpha-\alpha_*|]$.
The verified condition $4v<\Dloc$ isolates the retained band and places
it below the chosen cutoff. Weyl bounds put every omitted state at least
$\Dloc-2v$ above the true ground energy. With the retained partition
function shifted by that energy,
\begin{equation}
 \frac{Z_{\rm out}}{Z_{\rm low}}
 \le \frac{(2^{5N}-4^N)e^{-(\Dloc-2v)/T}}{Z_{\rm low}}.
\end{equation}
The logarithm of this bound is below $-2167$ in all tested thermal cases;
the tail was bounded, not set to zero by numerical underflow.
The raw positive-frequency cutoff was $10^{-10}$ in $J_1$ units.
An audit at $10^{-11}$ and $10^{-9}$ exposed the quasielastic contribution
at $\tau=1.6$. The subsequent fixed scaled windows are reported separately
and do not replace the initial protocol or its data.

\section{Relation to the original full spectra}

The original $L=40$ maps in Fig.~\ref{fig:rank-original} have
$\beta/\Dloc\simeq1.46$ and $3.96$, outside the validated perturbative
range, and finite-time resolution broader than the separately calculated
open-chain gaps. They neither establish the weak-coupling powers in
Fig.~\ref{fig:rank-evidence} nor resolve the tiny retuned weights; they
are retained as context, not validation of the composite-probe expansion.


\bibliography{references,rank_references}
\end{document}
